
\section{Spin angular-momentum density: the axial component}
\label{sec:canonical-sam-density-and-flux}

The spin angular-momentum (SAM) density in the time domain is
\begin{align}
  \boldsymbol{\mathcal S}=\epsilon_0{\bf E}\times{\bf A},\qquad
  \mathcal S_z=\epsilon_0(E_x A_y-E_y A_x).
  \label{eq:canonical-sam-time-domain}
\end{align}
We use the radiation gauge introduced in
Eq.~\eqref{eq:am-general-radiation-gauge}.  Tildes denote Fourier amplitudes
at $(\omega,{\bf x})$, with $\omega>0$, and
$\tilde{\bf A}=\tilde{\bf E}/(i\omega)$.
In spectral expressions, $\mathcal S_z$ denotes the time-integrated density.
With the Fourier convention of
Eq.~\eqref{eq:Faraday-Fourier-convention-observables}, the one-sided
bilinear spectral factor is $4\pi$, giving
\begin{align}
  \frac{d\mathcal S_z}{d\omega}
  =4\pi\epsilon_0\operatorname{Re}\left[
    \tilde E_x\tilde A_y^*-\tilde E_y\tilde A_x^*
  \right].
  \label{eq:canonical-sam-bilinear-spectrum}
\end{align}
Substituting $\tilde A_i^*=i\tilde E_i^*/\omega$ yields
\begin{align}
  \tilde E_x\tilde A_y^*-\tilde E_y\tilde A_x^*
  &=\frac{i}{\omega}\left(
    \tilde E_x\tilde E_y^*-\tilde E_y\tilde E_x^*\right)\nonumber\\
  &=-\frac{2}{\omega}\operatorname{Im}
    \left[\tilde E_x\tilde E_y^*\right]
  =\frac{2}{\omega}\operatorname{Im}
    \left[\tilde E_x^*\tilde E_y\right].
  \label{eq:canonical-sam-spectral-reduction}
\end{align}
Thus the axial SAM spectrum is
\begin{importantbox}
  \begin{align}
    \frac{d\mathcal S_z}{d\omega}
    =\frac{8\pi\epsilon_0}{\omega}\operatorname{Im}
      \left[\tilde E_x^*\tilde E_y\right]
    =-\frac{8\pi\epsilon_0}{\omega}\operatorname{Im}
      \left[\tilde E_x\tilde E_y^*\right],\qquad \omega>0.
    \label{eq:canonical-sam-density-spectrum}
  \end{align}
\end{importantbox}

\section{Spin angular-momentum flux and continuity equation}
\label{sec:canonical-sam-flux}

The local conservation law for spin angular momentum in vacuum is
\begin{importantbox}
  \begin{align}
    \partial_t\mathcal S_i+\partial_j\Sigma_{ij}=0.
    \label{eq:canonical-sam-continuity}
  \end{align}
\end{importantbox}
Here $\mathcal S_i=\epsilon_0({\bf E}\times{\bf A})_i$ is the instantaneous
spin density, and $\Sigma_{ij}$ is the flux of its $i$th component in the
$j$th direction.  The spin flux tensor in SI units is
\begin{align}
  \Sigma_{ij}
  &=\epsilon_0c^2\left[
    \delta_{ij}({\bf B}\cdot{\bf A})-B_i A_j-B_j A_i\right]\nonumber\\
  &=\frac{1}{\mu_0}\left[
    \delta_{ij}({\bf B}\cdot{\bf A})-B_i A_j-B_j A_i\right].
  \label{eq:canonical-sam-flux-tensor}
\end{align}
Repeated spatial indices are summed.

\subsection{Derivation of the continuity equation}

Differentiating the spin density in the radiation gauge gives
\begin{align}
  \partial_t\boldsymbol{\mathcal S}
  =\epsilon_0\left[
    (\partial_t{\bf E})\times{\bf A}
    +{\bf E}\times(\partial_t{\bf A})\right].
\end{align}
The second term vanishes because $\partial_t{\bf A}=-{\bf E}$.
Using the vacuum Maxwell--Amp\`ere equation,
$\partial_t{\bf E}=c^2\boldsymbol{\nabla}\times{\bf B}$, we obtain
\begin{align}
  \partial_t\mathcal S_i
  &=\epsilon_0c^2\left[
    (\boldsymbol{\nabla}\times{\bf B})\times{\bf A}\right]_i\nonumber\\
  &=\epsilon_0c^2\left[
    (\partial_l B_i)A_l-(\partial_i B_l)A_l\right].
  \label{eq:canonical-sam-time-derivative}
\end{align}
The divergence of the flux tensor is
\begin{align}
  \partial_j\Sigma_{ij}
  &=\epsilon_0c^2\partial_j\left[
    \delta_{ij}B_l A_l-B_i A_j-B_j A_i\right]\nonumber\\
  &=\epsilon_0c^2\left[
    (\partial_i B_l)A_l-(\partial_l B_i)A_l
    +B_j(\partial_i A_j-\partial_j A_i)\right],
\end{align}
where $\partial_j A_j=0$ and $\partial_j B_j=0$ were used.
The remaining term containing derivatives of the potential vanishes:
\begin{align}
  B_j(\partial_i A_j-\partial_j A_i)
  =\left[{\bf B}\times(\boldsymbol{\nabla}\times{\bf A})\right]_i
  =({\bf B}\times{\bf B})_i=0.
\end{align}
Consequently,
\begin{highlightbox}
  \begin{align}
    \partial_j\Sigma_{ij}
    =\epsilon_0c^2\left[
      (\partial_i B_l)A_l-(\partial_l B_i)A_l\right]
    =-\partial_t\mathcal S_i,
    \label{eq:canonical-sam-continuity-verification}
  \end{align}
\end{highlightbox}
which establishes Eq.~\eqref{eq:canonical-sam-continuity}.

\subsection{Spectral density of the axial flux}

The flux along $Oz$ of the spin component $\mathcal S_z$ is obtained by
setting $i=j=z$:
\begin{highlightbox}
  \begin{align}
    \Sigma_{zz}
    &=\epsilon_0c^2\left[({\bf B}\cdot{\bf A})-2B_z A_z\right]\nonumber\\
    &=\epsilon_0c^2\left[B_x A_x+B_y A_y-B_z A_z\right].
    \label{eq:canonical-sam-axial-flux}
  \end{align}
\end{highlightbox}
As for the density, $d\Sigma_{zz}/d\omega$ denotes the one-sided spectrum
of the time-integrated flux.  Applying the bilinear spectral rule gives
\begin{align}
  \frac{d\Sigma_{zz}}{d\omega}
  =4\pi\epsilon_0c^2\operatorname{Re}\left[
    \tilde B_x\tilde A_x^*+\tilde B_y\tilde A_y^*
    -\tilde B_z\tilde A_z^*\right].
  \label{eq:canonical-sam-flux-bilinear}
\end{align}
Substituting $\tilde A_k^*=i\tilde E_k^*/\omega$, the result is
\begin{importantbox}
  \begin{align}
    \frac{d\Sigma_{zz}}{d\omega}
    &=-\frac{4\pi\epsilon_0c^2}{\omega}\operatorname{Im}\left[
      \tilde B_x\tilde E_x^*+\tilde B_y\tilde E_y^*
      -\tilde B_z\tilde E_z^*\right]\nonumber\\
    &=\frac{4\pi\epsilon_0c^2}{\omega}\operatorname{Im}\left[
      \tilde B_x^*\tilde E_x+\tilde B_y^*\tilde E_y
      -\tilde B_z^*\tilde E_z\right],\qquad \omega>0.
    \label{eq:canonical-sam-flux-spectrum}
  \end{align}
\end{importantbox}
All Fourier amplitudes are evaluated at $(\omega,{\bf x})$.
